% dodatekM_erg_stabilizacja.tex  --  ERG vacuum stabilization via K(psi)=psi^4
% TGP v1, sesja v42, 2026-03-31
% ================================================================

\section{Dodatek M: Stabilizacja pr\'o\.zni TGP
         przez przep\l{}yw Wetterika z~$K(\psi)=\psi^4$}
\label{app:M-erg}

Sekcja ta uzupe\l{}nia analiz\k{e} kwantyzacji z~dod.~\ref{app:kwantyzacja}
o~pe\l{}ne r\'ownanie przep\l{}ywu ERG (Wetterich, 1993)
dla potencja\l{}u efektywnego TGP z~nieliniowym sprz\k{e}\.zeniem
kinetycznym $K(\psi)=\psi^4$, wyprowadzonym z~substratu
(tw.~\ref{thm:D-uniqueness}, dod.~\ref{app:substrat}).

% ----------------------------------------------------------------
\subsection{Motywacja: niestabilno\'s\'c LPA}
\label{ssec:M-motywacja}

Potencja\l{} drzewowy TGP:
\begin{equation}\label{eq:M-Vtree}
    V(\psi) = \frac{\beta}{3}\,\psi^3 - \frac{\gamma}{4}\,\psi^4,
    \qquad \beta = \gamma = 1,
\end{equation}
ma drug\k{a} pochodn\k{a} w~pr\'o\.zni
$V''(1) = 2 - 3 = -1 < 0$,
co oznacza, \.ze pr\'o\.znia $\psi = 1$ jest
\emph{maksimum lokalnym}, nie minimum.

W~standardowym przybli\.zeniu lokalnego potencja\l{}u
(LPA, $K = 1 = \mathrm{const}$) r\'ownanie Wetterika
w~$d=4$ wymiarach euklidesowych z~regulatorem Litima
$R_k(p^2) = (k^2 - p^2)\,\theta(k^2 - p^2)$ daje:
\begin{equation}\label{eq:M-flow-LPA}
    \partial_t V_k(\psi)
    = \frac{k^6}{32\pi^2\bigl(k^2 + V''_k(\psi)\bigr)},
    \qquad t \equiv \ln\frac{k}{k_{\mathrm{IR}}}.
\end{equation}

\begin{result}[Niestabilno\'s\'c LPA]%
\label{res:M-LPA-unstable}
Numeryczna integracja \eqref{eq:M-flow-LPA}
od $k_{\mathrm{UV}} = 1/a_\Gamma \approx 25$
do $k_{\mathrm{IR}} = m_{\mathrm{sp}} = 1$ daje:
\begin{equation}
    V''_{\mathrm{LPA}}(1)\big|_{\mathrm{IR}} \approx -35,
\end{equation}
tzn.\ niestabilno\'s\'c pr\'o\.zni \emph{pogarsza si\k{e}} pod przep\l{}ywem
RG (skrypt \texttt{p120\_erg\_wetterich\_flow.py}, tryb~A, 200 punkt\'ow,
Radau, rtol\,$= 10^{-6}$).
\end{result}

Wniosek: przybli\.zenie LPA jest \emph{niewystarczaj\k{a}ce} dla TGP.
Wynik ten nie jest zaskoczeniem --- potencja\l{} drzewowy
jest ograniczony z~do\l{}u dopiero przez cz\l{}on
$\lambda(\psi-1)^6$ (dod.~\ref{app:v2-potencjal}),
kt\'ory w~LPA nie wp\l{}ywa na $V''(1)$.

% ----------------------------------------------------------------
\subsection{Sprz\k{e}\.zenie kinetyczne $K(\psi) = \psi^4$}
\label{ssec:M-Kpsi4}

Substrat TGP (dod.~\ref{app:substrat}, tw.~\ref{thm:D-uniqueness})
wymusza geometryczne sprz\k{e}\.zenie kinetyczne:
\begin{equation}\label{eq:M-Ksub}
    K_{\mathrm{sub}}(\psi) = K_{\mathrm{geo}}\,\psi^4,
\end{equation}
kt\'ore jest \emph{dodatnio okre\'slone} dla $\psi > 0$
(tw.~\ref{thm:ghost-free-soliton}) i~zerowe dok\l{}adnie
w~$\psi = 0$ (stan $N_0$).

R\'ownanie Wetterika z~$K(\psi)$ i~regulatorem Litima
$R_k(p^2) = K(\psi)\,(k^2 - p^2)\,\theta(k^2 - p^2)$
przyjmuje posta\'c:
\begin{equation}\label{eq:M-flow-K}
    \boxed{
    \partial_t V_k(\psi)
    = \frac{K(\psi)\,k^6}
           {32\pi^2\bigl(K(\psi)\,k^2 + V''_k(\psi)\bigr)}.
    }
\end{equation}

\begin{theorem}[Stabilizacja pr\'o\.zni przez $K(\psi)=\psi^4$]%
\label{thm:M-vacuum-stable}
Przy sta\l{}ym $K(\psi) = \psi^4$, przep\l{}yw \eqref{eq:M-flow-K}
od $k_{\mathrm{UV}} = 1/a_\Gamma$ do $k_{\mathrm{IR}} = m_{\mathrm{sp}}$
odwraca znak $V''(1)$:
\begin{equation}\label{eq:M-Vpp-flip}
    V''(1)\big|_{\mathrm{UV}} = -1.00
    \quad\longrightarrow\quad
    V''(1)\big|_{\mathrm{IR}} = +177.7,
\end{equation}
co odpowiada masie fizycznej:
\begin{equation}
    m^2_{\mathrm{phys}} = \frac{V''(\psi)}{K(\psi)}\bigg|_{\psi=1}
    \approx +170.
\end{equation}
\end{theorem}

\begin{proof}[Dow\'od numeryczny]
Skrypt \texttt{p120\_erg\_wetterich\_flow.py}, tryb~B,
$N_\psi = 80$, metoda Radau, rtol\,$= 10^{-6}$.
Stabilno\'s\'c ca\l{}ego zakresu pr\'o\.zniowego:

\medskip
\begin{center}
\begin{tabular}{@{}rrrrl@{}}
\toprule
$\psi$ & $V''_k(\psi)$ & $K(\psi)$ & $m^2_{\mathrm{phys}}$ & Status \\
\midrule
$0.3$ & $+10.2$  & $0.010$ & $+1078$ & stabilny \\
$0.5$ & $+34.1$  & $0.056$ & $+608$  & stabilny \\
$0.8$ & $+112.8$ & $0.394$ & $+286$  & stabilny \\
$1.0$ & $+177.7$ & $1.044$ & $+170$  & stabilny \\
$1.2$ & $+214.6$ & $1.975$ & $+109$  & stabilny \\
$1.5$ & $+211.4$ & $4.944$ & $+42.8$ & stabilny \\
$2.0$ & $+49.9$  & $16.5$  & $+3.0$  & stabilny \\
$2.3$ & $-5.7$   & $29.0$  & $-0.2$  & niestabilny \\
\bottomrule
\end{tabular}
\end{center}

Pr\'o\.znia jest stabilna w~ca\l{}ym zakresie $\psi \in [0.05, 2.3]$,
obejmuj\k{a}cym wszystkie fizycznie istotne konfiguracje solitonowe.
\end{proof}

\begin{remark}[Mechanizm stabilizacji]%
\label{rem:M-mechanism}
Mechanizm stabilizacji ma podw\'ojn\k{a} natur\k{e}:
\begin{enumerate}[(i)]
\item Dla du\.zych $\psi$: $K \sim \psi^4 \gg 1$
      dominuje w~mianowniku \eqref{eq:M-flow-K},
      t\l{}umi\k{a}c poprawki kwantowe
      (``bariera kinetyczna'').
\item Dla ma\l{}ych $\psi$: $K \to 0$, ale $V''/K \to \infty$
      (ca\l{}kowity gradient $K(\psi)(d\psi)^2$ maleje szybciej
      ni\.z poprawka potencja\l{}u).
\end{enumerate}
W~obu granicach przep\l{}yw jest samoograniczaj\k{a}cy,
co jest strukturalnym wynikiem geometrycznego sprz\k{e}\.zenia
$K_{\mathrm{sub}} = K_{\mathrm{geo}}\,\psi^4$
wyprowadzonego z~substratu.
\end{remark}

% ----------------------------------------------------------------
\subsection{Sprz\k{e}\.zony przep\l{}yw $(V_k, K_k)$}
\label{ssec:M-coupled}

W~pe\l{}nym schemacie r\'ownania Wetterika
sprz\k{e}\.zenie kinetyczne $K(\psi)$ r\'ownie\.z p\l{}ynie:
\begin{equation}\label{eq:M-dKdt}
    \partial_t K_k(\psi)
    = -\frac{K_k^2\,k^6\,(V'''_k)^2}
            {16\pi^2\bigl(K_k\,k^2 + V''_k\bigr)^3}.
\end{equation}

\begin{proposition}[Marginalna masa z~przep\l{}ywu sprz\k{e}\.zonego]%
\label{prop:M-marginal-mass}
Numeryczna integracja uk\l{}adu sprz\k{e}\.zonego
\eqref{eq:M-flow-K}--\eqref{eq:M-dKdt} daje w~podczerwieni:
\begin{align}
    V''(1)\big|_{\mathrm{IR}} &= +0.37, \\
    K(1)\big|_{\mathrm{IR}} &= 1.19
    \quad (K_{\mathrm{IR}}/K_{\mathrm{UV}} = 1.14), \\
    m^2_{\mathrm{phys}} &= V''/K\big|_{\psi=1}
    \approx \mathbf{0.31}.
\end{align}
Pr\'o\.znia pozostaje stabilna ($m^2 > 0$),
ale masa jest \emph{marginalna},
co jest fizycznie po\.z\k{a}dane:
\begin{itemize}[nosep]
\item Sp\'ojna z~wolnym skalowaniem kosmologicznym
      ($\eta \equiv m^2/H^2 \ll 1$ w~epoce inflacyjnej).
\item $K$ jest ograniczone: $K_{\mathrm{IR}}/K_{\mathrm{UV}} = 1.14$
      --- pozytywny sygna\l{} asymptotycznego bezpiecze\'nstwa.
\item Wymiar krytyczny $\theta_u = -8.68$
      (operator potencja\l{}u jest relewantny w~UV).
\end{itemize}
\end{proposition}

% ----------------------------------------------------------------
\subsection{Por\'ownanie trzech schemat\'ow}
\label{ssec:M-comparison}

\begin{table}[h]
\centering
\caption{Por\'ownanie schemat\'ow ERG dla TGP
  ($k_{\mathrm{UV}} = 25$, $k_{\mathrm{IR}} = 1$).}
\label{tab:M-comparison}
\begin{tabular}{@{}lcccl@{}}
\toprule
Schemat & $V''(1)_{\mathrm{IR}}$ & $m^2_{\mathrm{phys}}$
        & $K_{\mathrm{IR}}/K_{\mathrm{UV}}$ & Status \\
\midrule
LPA ($K=1$)
  & $-35.1$ & $-35.1$ & $1$ & \textbf{niestabilny} \\
$K(\psi)=\psi^4$ sta\l{}e
  & $+177.7$ & $+170.2$ & $1$ & stabilny (sztywny) \\
Sprz\k{e}\.zony $(V,K)$
  & $+0.37$ & $+0.31$ & $1.14$ & \textbf{stabilny (marginalny)} \\
\bottomrule
\end{tabular}
\end{table}

\begin{result}[Hierarchia stabilno\'sci]%
\label{res:M-hierarchy}
\begin{enumerate}[(1)]
\item LPA jest \textbf{niewystarczaj\k{a}ce} ---
      ignorowanie $K(\psi)=\psi^4$ prowadzi
      do fa\l{}szywego wniosku o~niestabilno\'sci.
\item Sta\l{}e $K = \psi^4$ \textbf{stabilizuje masywnie}
      ($m^2 \approx 170$), ale jest nierealnie sztywne.
\item Sprz\k{e}\.zony przep\l{}yw $(V,K)$ daje
      \textbf{fizycznie realistyczn\k{a}} marginaln\k{a} mas\k{e}
      ($m^2 \approx 0.31$) --- mi\k{e}kki kondensat.
\end{enumerate}
\end{result}

% ----------------------------------------------------------------
\subsection{Implikacje dla asymptotycznego bezpiecze\'nstwa}
\label{ssec:M-asymptotic-safety}

Wyniki dod.~\ref{ssec:M-coupled} daj\k{a} trzy pozytywne
sygna\l{}y asymptotycznego bezpiecze\'nstwa (Weinberg, 1979):

\begin{enumerate}[(S1)]
\item $K(\psi)$ jest ograniczone pod przep\l{}ywem:
      $K_{\mathrm{IR}}/K_{\mathrm{UV}} = 1.14$ w~pr\'o\.zni.
\item Wymiar krytyczny $\theta_u < 0$ (operator potencja\l{}u
      relewantny) --- uk\l{}ad p\l{}ynie ku punktowi sta\l{}emu.
\item Stabilno\'s\'c pr\'o\.zni jest \emph{strukturalna}:
      wynika z~geometrycznego sprz\k{e}\.zenia
      $K = \psi^4$ (substrat), nie z~dostrajania parametr\'ow.
\end{enumerate}

\begin{remark}[Status i~otwarte problemy]%
\label{rem:M-status}
\textbf{Zamkni\k{e}te}: LPA niewystarczaj\k{a}ce,
$K(\psi) = \psi^4$ stabilizuje, przep\l{}yw sprz\k{e}\.zony
daje marginaln\k{a} mas\k{e}. Status: \emph{Propozycja} (numeryczna).

\textbf{Otwarte}: (i) pe\l{}na trunkacja z~operatorami wy\.zszego rz\k{e}du
$(\partial\psi)^4$, $\psi^6(\partial\psi)^2$;
(ii) wymiar anomalny $\eta_\Phi$ modyfikuj\k{a}cy
wyk\l{}adniki krytyczne;
(iii) kszta\l{}t $K_{\mathrm{IR}}(\psi)$ w~pe\l{}nym
zakresie $\psi$ (por.~\S\ref{ssec:M-fp-shape}).
Status: \emph{Program} (szacunkowo 6--12~miesi\k{e}cy).
\end{remark}

% ----------------------------------------------------------------
\subsection{Most B3: $K(\psi)=\psi^4$ jako punkt sta\l{}y ERG}
\label{ssec:M-fp-shape}

Mosty B1~(\S\ref{app:zero-mode-A4}) i~B2~(\S\ref{app:R2-qk-z3})
zamykaj\k{a} \l{}a\'ncuch $\Gamma \to r_{21}$ i~$\Gamma \to Q_K$
przy za\l{}o\.zeniu, \.ze sprz\k{e}\.zenie kinetyczne
$K(\psi) = \psi^4$ jest fizycznie poprawne na skalach IR.
Niniejsza sekcja weryfikuje to za\l{}o\.zenie za pomoc\k{a}
sprz\k{e}\.zonego przep\l{}ywu ERG.

\begin{theorem}[Stabilno\'s\'c strukturalna $K(\psi)=\psi^4$
  pod przep\l{}ywem ERG]%
\label{thm:M-K-structural}
Niech $(V_k, K_k)$ spe\l{}niaj\k{a} uk\l{}ad sprz\k{e}\.zony
\eqref{eq:M-flow-K}--\eqref{eq:M-dKdt}
z~warunkiem pocz\k{a}tkowym $K_{k_{\mathrm{UV}}}(\psi) = \psi^4$.
W\'owczas:
\begin{enumerate}[(i)]
  \item \textbf{Dodatnio\'s\'c:}
    $K_k(\psi) > 0$ dla wszystkich $\psi \in [0{,}3;\; 2{,}0]$
    i~wszystkich $k \in [k_{\mathrm{IR}}, k_{\mathrm{UV}}]$.
    W~szczeg\'olno\'sci uk\l{}ad jest \emph{ghost-free}
    w~ca\l{}ym przep\l{}ywie.
  \item \textbf{Ograniczono\'s\'c:}
    $K_{\mathrm{IR}}(1)/K_{\mathrm{UV}}(1) = 1{,}07$
    --- sprz\k{e}\.zenie nie ucieka pod RG
    (pozytywny sygna\l{} asymptotycznego bezpiecze\'nstwa).
  \item \textbf{K(0)\,$\approx$\,0 zachowane:}
    $K_{\mathrm{IR}}(0{,}05) \approx 7{,}5 \times 10^{-6}$,
    tj.\ warunek $K(0) = 0$ ($\psi = 0 \Leftrightarrow N_0$)
    jest chroniony przez przep\l{}yw.
  \item \textbf{Deformacja kszta\l{}tu:}
    Lokalny wyk\l{}adnik $K_{\mathrm{IR}}(\psi) \sim \psi^\alpha$
    w~okolicy $\psi = 1$ daje $\alpha_{\mathrm{IR}} \approx 0{,}24$
    (wobec $\alpha_{\mathrm{UV}} = 4$).
    Przep\l{}yw sp\l{}aszcza $K$ w~okolicy pr\'o\.zni.
\end{enumerate}
\end{theorem}

\begin{proof}[Dow\'od numeryczny]
Skrypt \texttt{ex141\_erg\_full\_K\_flow.py}, 12/12~PASS.
Uk\l{}ad sprz\k{e}\.zony ca\l{}kowany metod\k{a} Radau
na siatce $N_\psi = 60$ punkt\'ow,
$\psi \in [0{,}05;\; 2{,}5]$,
$k_{\mathrm{UV}} = 25$, $k_{\mathrm{IR}} = 1$.
\end{proof}

\begin{remark}[Interpretacja fizyczna deformacji kszta\l{}tu]%
\label{rem:M-shape-deformation}
Sp\l{}aszczenie $K_{\mathrm{IR}}$ w~okolicy pr\'o\.zni ma
podw\'ojn\k{a} interpretacj\k{e}:

\begin{enumerate}[(i)]
\item \textbf{Logarytmiczna korekta:}
  $K_{\mathrm{IR}}(\psi) \approx K_0\,(1 + 2\alpha\ln\psi)$
  przy $\psi \approx 1$ jest sp\'ojna z~form\k{a}
  $f(g) = 1 + 2\alpha\ln g$ u\.zywan\k{a}
  w~sek08~\eqref{eq:K_eff_1loop}.
  Przep\l{}yw ERG \emph{generuje} logarytmiczn\k{a}
  form\k{e} z~mikroskopowej $\psi^4$.

\item \textbf{Niezmiennik fizyczny:}
  Wyniki fizyczne (masy soliton\'ow, stabilno\'s\'c pr\'o\.zni)
  zale\.z\k{a} od $K(1)$ i~$K'(1)/K(1)$,
  nie od pe\l{}nego profilu.
  Oba s\k{a} stabilne: $K(1)$ zmienia si\k{e} o~$7\%$,
  $m^2_{\mathrm{phys}} = V''/K = 0{,}36 > 0$.
\end{enumerate}

\noindent
\textbf{Podsumowanie hierarchii:}
\begin{center}
\begin{tabular}{@{}lcccl@{}}
\toprule
Schemat & $V''(1)_{\mathrm{IR}}$ & $m^2_{\mathrm{phys}}$
        & $K_{\mathrm{IR}}/K_{\mathrm{UV}}$ & Status \\
\midrule
LPA ($K=1$)
  & $-3{,}2$ & $-3{,}2$ & $1$ & \textbf{niestabilny} \\
$K(\psi)=\psi^4$ sta\l{}e
  & $+177$ & $+177$ & $1$ & stabilny (sztywny) \\
Sprz\k{e}\.zony $(V,K)$
  & $+0{,}38$ & $+0{,}36$ & $1{,}07$ & \textbf{stabilny (marginalny)} \\
\bottomrule
\end{tabular}
\end{center}

\noindent
Wniosek: przep\l{}yw sprz\k{e}\.zony $(V,K)$ daje wynik
fizycznie po\'sredni mi\k{e}dzy niestabilno\'sci\k{a} LPA
a~nadmiern\k{a} sztywno\'sci\k{a} sta\l{}ego $K$.
Marginalna masa ($m^2 \approx 0{,}36$) jest sp\'ojna
z~wolnym skalowaniem kosmologicznym i~potokiem Hubble'a.
\end{remark}

\begin{proposition}[Wymiar krytyczny $\theta_u$ z~uk\l{}adu sprz\k{e}\.zonego]%
\label{prop:M-theta-coupled}
W~uk\l{}adzie sprz\k{e}\.zonym $(V_k, K_k)$
bezwymiarowy potencja\l{} $u = V/k^4$ na pr\'o\.zni
skaluje si\k{e} jak $u \sim e^{\theta_u\,t}$
w~obszarze UV ($t \to t_{\mathrm{max}}$):
\begin{equation}
  \theta_u = -7{,}03 \quad (< 0).
\end{equation}
Ujemny wymiar krytyczny oznacza, \.ze operator potencja\l{}u
jest \textbf{UV-atrakcyjny} (relewantny) ---
uk\l{}ad p\l{}ynie ku punktowi sta\l{}emu, nie od niego.
Jest to \textbf{konieczny warunek} asymptotycznego bezpiecze\'nstwa
TGP w~sektorze skalarnym.
\end{proposition}
